% META: {"id": "TCOMPL-AIR-EX-015", "chapitre": "TCOMPL-CALCULS-AIRES", "type_objet": "exercice", "capacites": ["TCOMPL-AIR-C3"], "capacites_codes": ["C3"], "methodes": ["M2","M5"], "parcours": 1, "competences": ["raisonner","calculer"], "duree_min": 18, "mode_creation": "ex_nihilo", "sources_inspiration": [], "fichier_tex": "chapitres/TCOMPL-CALCULS-AIRES/exercices/TCOMPL-AIR-EX-015.tex", "corrige_tex": "chapitres/TCOMPL-CALCULS-AIRES/corriges/TCOMPL-AIR-CO-015.tex", "status": "approved"}
% BEGIN-VERIFY
% from sympy import *
% x = symbols('x', positive=True)
% f = ln(x)/x
% fp = diff(f, x)
% assert simplify(fp) == (1 - ln(x))/x**2
% assert solve(fp, x) == [exp(1)]
% assert f.subs(x, exp(1)) == exp(-1)
% assert limit(f, x, oo) == 0
% assert limit(f, x, 0, '+') == -oo
% END-VERIFY
\begin{exercice}{TCOMPL-AIR-EX-015}{2}{18}
Soit $f(x) = \dfrac{\ln x}{x}$ sur $\left]0, +\infty\right[$.
\begin{enumerate}
  \item Calculer $f'(x)$ et etudier son signe.
  \item Determiner les limites de $f$ en $0^+$ et en $+\infty$.
  \item Dresser le tableau de variations.
\end{enumerate}

\end{exercice}
